% Generated draft: requires mathematical and source review.
\documentclass[11pt]{article}
\usepackage[margin=25mm]{geometry}
\usepackage{fontspec}
\setmainfont{Latin Modern Roman}
\usepackage{amsmath,amssymb,mathtools,unicode-math}
\setmathfont{Latin Modern Math}
\usepackage{enumitem,needspace,xurl,hyperref}
\setlength{\parindent}{0pt}
\setlength{\parskip}{5pt}
\setlength{\emergencystretch}{3em}
\newcommand{\N}{\mathbb{N}}
\newcommand{\Z}{\mathbb{Z}}
\newcommand{\Q}{\mathbb{Q}}
\newcommand{\R}{\mathbb{R}}
\newcommand{\C}{\mathbb{C}}
\newcommand{\F}{\mathbb{F}}
\newcommand{\A}{\mathbb{A}}
\newcommand{\PP}{\mathbb{P}}
\newcommand{\E}{\mathbb{E}}
\newcommand{\eps}{\varepsilon}
\newcommand{\dd}{\,\mathrm{d}}
\begin{document}
\section*{UnsolvedMath ID: 701771}
\textbf{Source problem code:} catalogue entry
\subsection{Definitions and mathematical statement}
The source uses \(c\) both as a generator and as the nilpotency-class parameter. Here the generator is renamed \(u\), and the parameter is renamed \(r\). Fix the commutator convention \([x,y]=x^{-1}y^{-1}xy\).

The following makes the presentation precise under the usual universal-quotient interpretation of its exponent and nilpotency conditions. For \((q,r)\in\{(8,12),(7,13)\}\), let \(F\) be the free group on \(a,b,u,d\). Define \(\gamma_1(F)=F\) and \(\gamma_{j+1}(F)=[\gamma_j(F),F]\), and let \(F^q\) be the normal subgroup generated by all \(q\)-th powers of elements of \(F\). Put
\[
H(q,r)=F\big/\Bigl(F^q\,\gamma_{r+1}(F)\,\langle\!\langle[b,a][d,u]^{-1}\rangle\!\rangle_F\Bigr),
\]
where \(\langle\!\langle w\rangle\!\rangle_F\) denotes the normal closure of \(w\) in \(F\). Thus this quotient is generated by the images of \(a,b,u,d\), satisfies \([b,a]=[d,u]\), has exponent dividing \(q\), and has nilpotency class at most \(r\).

The source says ``of exponent \(q\)'' and ``nilpotent of class \(c\)''. It does not explicitly distinguish the universal quotient above from requirements of exact exponent and exact nilpotency class. Those exactness requirements, and the intended universal interpretation, require confirmation; they are not established here.

For a group \(G\), its Schur multiplier is \(M(G)=H_2(G,\Z)\), with trivial integral coefficients. For an abelian group \(A\), its exponent is the least positive integer \(n\), if one exists, such that \(nx=0\) for every \(x\in A\).

Under the stated interpretation of \(H(q,r)\), the two questions attributed to M. R. Vaughan-Lee are:
\begin{enumerate}
\item[(a)] Is \(\operatorname{exp} M(H(8,12))=32\)?
\item[(b)] Is \(\operatorname{exp} M(H(7,13))=49\)?
\end{enumerate}
\subsection{Short English statement}
Vaughan-Lee asks whether the Schur multipliers of the two specified four-generator groups have exponents exactly \(32\) and \(49\), respectively. The draft uses an explicitly qualified universal-quotient interpretation of the presentation; the intended exponent and nilpotency-class conventions require review.
\subsection{Sources}
\begin{itemize}
\item Local UnsolvedMath registry, supplied complete source record: permanent ID \(701771\), exact source code \texttt{catalogue entry}, title \emph{Kourovka Notebook Problem 20.113}; upstream revision \texttt{d371ee917ac9647fbb06cd6fc529deec6f3dda82}. The record supplies the status label \texttt{open}; this is not an independently verified assessment.
\item M. R. Vaughan-Lee, proposer, 2022. The supplied attribution identifies \emph{The Kourovka Notebook (Unsolved Problems in Group Theory)}, 20th issue (2022), Problem 20.113, and the 21st issue edition (2026), identified as \texttt{arXiv:1401.0300v48}. These bibliographic details are reproduced from the record without independent source consultation.
\item The supplied statement reports that the groups were too large for the versions of the \(p\)-Quotient Algorithm described there as using 32-bit arithmetic, and suggests that a version using 64-bit arithmetic would help. This is a source-specific computational observation, not a verified assessment of present software capabilities. The registry record was updated on \texttt{2026-10-16}; that timestamp is not an independent literature-check date.
\end{itemize}
\end{document}
